% Source: physical PDF pp. 91--98 (printed pp. 68--75), from the Section 26.3
% heading through the final paragraph immediately before Section 26.4.
% Inventory: Eqs. (26.3.1)--(26.3.45), all unnumbered displays, one footnote,
% one centered three-asterisk divider, and no figures, tables, or citations.
% Uncertainties: none.

\subsection{Chiral and Linear Superfields}

In the previous section we found that the presence of $D$ and $\lambda$
components in a general superfield stood in the way of using such superfields
in a physically satisfactory Lagrangian density. Suppose then that we consider
a superfield with
\begin{equation}
  \lambda=D=0\,.
  \tag{26.3.1}\label{eq:26.3.1}
\end{equation}
Are these conditions preserved by supersymmetry transformations? According to
Eqs.\linebreak \eqref{eq:26.2.17} and \eqref{eq:26.2.16}, the condition $D=0$ is
invariant if $\lambda=0$, but the condition $\lambda=0$ is only invariant if
we also impose the condition that
$\partial_\mu V_\nu-\partial_\nu V_\mu=0$, which requires that $V_\mu$ be a
pure gauge:
\begin{equation}
  V_\mu(x)=\partial_\mu Z(x)\,.
  \tag{26.3.2}\label{eq:26.3.2}
\end{equation}
Eq.~\eqref{eq:26.2.15} shows that, with $\lambda=0$, this condition is
preserved by a supersymmetry transformation. We have thus arrived at a reduced
superfield, subject to the constraints \eqref{eq:26.3.1} and
\eqref{eq:26.3.2}, with component fields having the transformation properties
\begin{equation}
  \delta C=i(\xi\omega-\bar\xi\bar\omega)\,,
  \tag{26.3.3}\label{eq:26.3.3}
\end{equation}
\begin{equation}
  \delta\omega_\alpha
  =(-M+iN)\xi_\alpha
   -(\sigma^\mu\bar\xi)_\alpha\partial_\mu(C+iZ)\,,
  \tag{26.3.4}\label{eq:26.3.4}
\end{equation}
\begin{equation}
  \delta M=i\left(
    \xi\sigma^\mu\partial_\mu\bar\omega
    +\bar\xi\bar\sigma^\mu\partial_\mu\omega\right)\,,
  \tag{26.3.5}\label{eq:26.3.5}
\end{equation}
\begin{equation}
  \delta N=\xi\sigma^\mu\partial_\mu\bar\omega
  -\bar\xi\bar\sigma^\mu\partial_\mu\omega\,,
  \tag{26.3.6}\label{eq:26.3.6}
\end{equation}
\begin{equation}
  \delta Z=\xi\omega+\bar\xi\bar\omega\,.
  \tag{26.3.7}\label{eq:26.3.7}
\end{equation}
Comparing with Eq.~\eqref{eq:26.1.21}, we see that this is the same as the
supermultiplet constructed by direct methods in Section 26.1, with the
identifications
\begin{equation}
  C=A,\qquad
  \omega_\alpha=-i\psi_\alpha,\qquad
  \bar\omega_{\dot\alpha}
    =i\widetilde{\bar\psi}_{\dot\alpha},\qquad
  M=G,\qquad
  N=-F,\qquad
  Z=B\,.
  \tag{26.3.8}\label{eq:26.3.8}
\end{equation}

A superfield satisfying the conditions \eqref{eq:26.3.1} and
\eqref{eq:26.3.2} is said to be
\emph{chiral}.\footnote{Some authors use the term `chiral' to describe a
special case of such superfields, introduced below, which are here called
left-chiral or right-chiral.  Here ``chiral'' denotes a superfield that can
be expressed as the sum of a left-chiral superfield and a right-chiral one.}

To distinguish a chiral superfield $X(x,\theta)$ from the general superfield
$S(x,\theta)$ of the previous section, we will use $A$, $B$, $F$, $G$, and
$\psi$ for its components, instead of $C$, $M$, $N$, $Z$, and $\omega$. By
using Eqs.~\eqref{eq:26.3.1}, \eqref{eq:26.3.2}, and \eqref{eq:26.3.8} in
Eq.~\eqref{eq:26.2.10}, we find the form of a general chiral superfield to be
\begin{equation}
  \begin{aligned}
  X(x,\theta,\bar\theta)={}&
  A-(\theta\psi+\bar\theta\widetilde{\bar\psi})
  +\frac12(\theta\theta+\bar\theta\bar\theta)F
  -\frac{i}{2}(\theta\theta-\bar\theta\bar\theta)G\\
  &+\theta\sigma_\mu\bar\theta\,\partial^\mu B
  +\frac{i}{2}\theta\theta\,
     \bar\theta\bar\sigma^\mu\partial_\mu\psi
  +\frac{i}{2}\bar\theta\bar\theta\,
     \theta\sigma^\mu\partial_\mu\widetilde{\bar\psi}\\
  &+\frac14\theta\theta\bar\theta\bar\theta\,\Box A\,.
  \end{aligned}
  \tag{26.3.9}\label{eq:26.3.9}
\end{equation}
Thus the restrictions \eqref{eq:26.3.1} remove the highest unconstrained
components, while \eqref{eq:26.3.2} turns the vector component into the
derivative term displayed above.
(We could just as well have taken $C=-B$, $\omega=\psi$, $M=-F$, $N=-G$,
and $Z=A$. We make the identifications \eqref{eq:26.3.8} because, as we see
here, for a scalar superfield they are consistent with the usual convention
that $A$ and $F$ are scalars while $B$ and $G$ are pseudoscalars.)

The chiral superfield \eqref{eq:26.3.9} may be further decomposed, as
\begin{equation}
  X(x,\theta,\bar\theta)=\frac{1}{\sqrt{2}}
  \left[\Phi(x,\theta,\bar\theta)
  +\widetilde{\Phi}(x,\theta,\bar\theta)\right]\,,
  \tag{26.3.10}\label{eq:26.3.10}
\end{equation}
where
\begin{equation}
  \begin{aligned}
  \Phi(x,\theta,\bar\theta)={}&
  \phi-\sqrt{2}\theta\psi+\theta\theta\,\mathcal F
  -i\theta\sigma^\mu\bar\theta\,\partial_\mu\phi\\
  &+\frac{i}{\sqrt{2}}\theta\theta\,
    \bar\theta\bar\sigma^\mu\partial_\mu\psi
  +\frac14\theta\theta\bar\theta\bar\theta\,\Box\phi\,,
  \end{aligned}
  \tag{26.3.11}\label{eq:26.3.11}
\end{equation}
\begin{equation}
  \begin{aligned}
  \widetilde{\Phi}(x,\theta,\bar\theta)={}&
  \widetilde{\phi}-\sqrt{2}\bar\theta\widetilde{\bar\psi}
  +\bar\theta\bar\theta\,\widetilde{\mathcal F}
  +i\theta\sigma^\mu\bar\theta\,\partial_\mu\widetilde{\phi}\\
  &+\frac{i}{\sqrt{2}}\bar\theta\bar\theta\,
    \theta\sigma^\mu\partial_\mu\widetilde{\bar\psi}
  +\frac14\theta\theta\bar\theta\bar\theta
    \,\Box\widetilde{\phi}\,,
  \end{aligned}
  \tag{26.3.12}\label{eq:26.3.12}
\end{equation}
with component fields defined by
\begin{equation}
  \phi\equiv\frac{A+iB}{\sqrt{2}},\qquad
  \psi_\alpha\ \hbox{transforms as }(1/2,0),\qquad
  \mathcal{F}\equiv\frac{F-iG}{\sqrt{2}}\,,
  \tag{26.3.13}\label{eq:26.3.13}
\end{equation}
\begin{equation}
  \widetilde{\phi}\equiv\frac{A-iB}{\sqrt{2}},\qquad
  \widetilde{\bar\psi}_{\dot\alpha}\ \hbox{transforms as }(0,1/2),
  \qquad
  \widetilde{\mathcal{F}}\equiv\frac{F+iG}{\sqrt{2}}\,.
  \tag{26.3.14}\label{eq:26.3.14}
\end{equation}
The component fields of either $\Phi$ or $\widetilde{\Phi}$ furnish complete
representations of the supersymmetry algebra:
\begin{equation}
  \delta\psi_\alpha
  =-i\sqrt{2}\,\sigma^\mu_{\alpha\dot\alpha}
    \bar\xi^{\dot\alpha}\partial_\mu\phi
  +\sqrt{2}\mathcal{F}\xi_\alpha\,,
  \tag{26.3.15}\label{eq:26.3.15}
\end{equation}
\begin{equation}
  \delta\mathcal{F}
  =-i\sqrt{2}\bar\xi\bar\sigma^\mu\partial_\mu\psi\,,
  \tag{26.3.16}\label{eq:26.3.16}
\end{equation}
\begin{equation}
  \delta\phi=\sqrt{2}\xi\psi\,,
  \tag{26.3.17}\label{eq:26.3.17}
\end{equation}
\begin{equation}
  \delta\widetilde{\bar\psi}_{\dot\alpha}
  =+i\sqrt{2}\,\xi^\alpha\sigma^\mu_{\alpha\dot\alpha}
    \partial_\mu\widetilde{\phi}
  +\sqrt{2}\widetilde{\mathcal{F}}\bar\xi_{\dot\alpha}\,,
  \tag{26.3.18}\label{eq:26.3.18}
\end{equation}
\begin{equation}
  \delta\widetilde{\mathcal{F}}
  =-i\sqrt{2}\xi\sigma^\mu
    \partial_\mu\widetilde{\bar\psi}\,,
  \tag{26.3.19}\label{eq:26.3.19}
\end{equation}
\begin{equation}
  \delta\widetilde{\phi}
  =\sqrt{2}\bar\xi\widetilde{\bar\psi}\,,
  \tag{26.3.20}\label{eq:26.3.20}
\end{equation}
Here the two independent index types display chirality directly:
\[
  \xi_\alpha:\ (\tfrac12,0),
  \qquad
  \bar\xi_{\dot\alpha}:\ (0,\tfrac12).
\]
Here the tilde distinguishes the independent fields in the right-chiral
multiplet; the dotted bar records their index type.  They become complex
conjugates of the left-chiral fields when the original superfield is real.
A superfield of the form \eqref{eq:26.3.11} or
\eqref{eq:26.3.12} is known as \emph{left-chiral} or \emph{right-chiral},
respectively. In the special case where a chiral superfield $X(x,\theta)$ is
real, its left-chiral and right-chiral parts $\Phi$ and $\widetilde{\Phi}$ are
complex conjugates, so that $\widetilde{\phi}=\phi^*$,
$\widetilde{\mathcal{F}}=\mathcal{F}^*$, and
\(\widetilde{\bar\psi}_{\dot\alpha}
=\bar\psi_{\dot\alpha}\equiv(\psi_\alpha)^*\).
However, if we do not require $X(x,\theta)$ to be real then in general there
is no relation between \(\Phi\) and \(\widetilde{\Phi}\), and
\(\psi_\alpha\) and \(\widetilde{\bar\psi}_{\dot\alpha}\) are independent;
it is even possible that one of the two superfields vanishes.

The component fields of the superfield $\Phi$ include two complex bosonic
components $\phi$ and $\mathcal{F}$, or four independent real bosonic
components, and one Weyl fermion field \(\psi_\alpha\), together with its
conjugate, which has four independent real fermionic components. This is
another example of the general result, derived
at the end of the previous section, that any set of fields that furnish a
representation of the supersymmetry algebra must have an equal number of
independent bosonic and fermionic components.

We can use Eqs.~\eqref{eq:26.A.5}, \eqref{eq:26.A.17}, and
\eqref{eq:26.A.18} to rewrite Eqs.~\eqref{eq:26.3.11} and
\eqref{eq:26.3.12} in a form that clarifies the way that these superfields
depend on \(\theta\) and \(\bar\theta\):
\begin{equation}
  \Phi(x,\theta,\bar\theta)
  =\phi(x_+)-\sqrt{2}\theta\psi(x_+)
  +\theta\theta\,\mathcal{F}(x_+)\,,
  \tag{26.3.21}\label{eq:26.3.21}
\end{equation}
\begin{equation}
  \widetilde{\Phi}(x,\theta,\bar\theta)
  =\widetilde{\phi}(x_-)
  -\sqrt{2}\bar\theta\widetilde{\bar\psi}(x_-)
  +\bar\theta\bar\theta\,\widetilde{\mathcal{F}}(x_-)\,,
  \tag{26.3.22}\label{eq:26.3.22}
\end{equation}
where
\begin{equation}
  x_\pm^\mu\equiv x^\mu
  \mp i\theta\sigma^\mu\bar\theta\,.
  \tag{26.3.23}\label{eq:26.3.23}
\end{equation}

The expansions of $\phi(x_+)$ and $\widetilde{\phi}(x_-)$ in powers of
$x^\mu-x_\pm^\mu$ terminate with quadratic terms, the expansions of
$\psi_\alpha(x_+)\) and
\(\widetilde{\bar\psi}_{\dot\alpha}(x_-)\) terminate with
linear terms, while the expansions of
$\mathcal{F}(x_+)$ and $\widetilde{\mathcal{F}}(x_-)$ terminate at zeroth
order, because all higher terms make contributions in
Eqs.~\eqref{eq:26.3.21} and \eqref{eq:26.3.22} that contain three or more
factors of \(\theta\) or \(\bar\theta\) and therefore vanish. For the same
reason, any superfield that depends only on \(\theta\) and \(x_+^\mu\), but
not otherwise on \(\bar\theta\), must take the form
\eqref{eq:26.3.21}, and any superfield that depends only on \(\bar\theta\)
and \(x_-^\mu\), but not otherwise on \(\theta\), must take the form
\eqref{eq:26.3.22}.

We have seen that for a superfield to be left-chiral or right-chiral is
entirely a matter of what the superfield is allowed to depend on. It follows
immediately that \emph{any function of left-chiral superfields (or of
right-chiral superfields), but not their complex conjugates or spacetime
derivatives, is a left-(or right-)chiral superfield.} This can also be shown in
a more formal way. Because \(\Phi\) depends on \(\bar\theta\) only through
\(x_+\), and \(\widetilde{\Phi}\) depends on \(\theta\) only through
\(x_-\), they satisfy
\begin{equation}
  \bar D_{\dot\alpha}\Phi=0,\qquad
  D_\alpha\widetilde{\Phi}=0\,,
  \tag{26.3.24}\label{eq:26.3.24}
\end{equation}
where
\begin{equation}
  \bar D_{\dot\alpha}
  =\frac{\partial}{\partial\bar\theta^{\dot\alpha}}
  -i\theta^\beta\sigma^\mu_{\beta\dot\alpha}\partial_\mu\,,
  \tag{26.3.25}\label{eq:26.3.25}
\end{equation}
\begin{equation}
  D_\alpha
  =-\frac{\partial}{\partial\theta^\alpha}
  +i\sigma^\mu_{\alpha\dot\beta}\bar\theta^{\dot\beta}\partial_\mu\,,
  \tag{26.3.26}\label{eq:26.3.26}
\end{equation}
for which
\[
  \bar D_{\dot\alpha}x_+^\mu
  =D_\alpha x_-^\mu=0\,.
\]
Conversely, if a superfield \(\Phi\) satisfies \(\bar D\Phi=0\) it is
left-chiral, and if it satisfies \(D\Phi=0\) it is right-chiral. Any
function $f(\Phi)$ of superfields $\Phi_n$ that all satisfy
\(\bar D\Phi_n=0\) or all satisfy \(D\Phi_n=0\) will satisfy
\(\bar Df(\Phi)=0\) or \(Df(\Phi)=0\), and hence be left-chiral or
right-chiral, respectively. But a function of left-chiral and right-chiral
superfields is in general not chiral at all.

Using the representation \eqref{eq:26.3.21} for left-chiral superfields makes
it easy to work out their multiplication properties. For instance, if
$\Phi_1$ and $\Phi_2$ are two left-chiral superfields, then their product
$\Phi=\Phi_1\Phi_2$ is a left-chiral superfield, with components
\begin{equation}
  \phi=\phi_1\phi_2\,,
  \tag{26.3.27}\label{eq:26.3.27}
\end{equation}
\begin{equation}
  \psi_\alpha=\phi_1\psi_{2\alpha}+\phi_2\psi_{1\alpha}\,,
  \tag{26.3.28}\label{eq:26.3.28}
\end{equation}
\begin{equation}
  \mathcal{F}
  =\phi_1\mathcal{F}_2+\phi_2\mathcal{F}_1
  -\psi_1\psi_2\,.
  \tag{26.3.29}\label{eq:26.3.29}
\end{equation}

The presence of chiral superfields in a theory opens up an additional
possibility for constructing supersymmetric actions. Inspection of the
transformation rule \eqref{eq:26.3.16} shows that a supersymmetry
transformation changes the $\mathcal{F}$-term of a left-chiral superfield
$\Phi$ by a derivative, so that the integral of the $\mathcal{F}$-term of any
left-chiral superfield is supersymmetric. Thus we can form a supersymmetric
action as
\begin{equation}
  I=\int d^4x\,[f]_{\mathcal{F}}
  +\int d^4x\,[f]_{\mathcal{F}}^*
  +\frac12\int d^4x\,[K]_D\,,
  \tag{26.3.30}\label{eq:26.3.30}
\end{equation}
where $f$ and $K$ are any left-chiral superfield and general real superfield,
respectively, formed from the elementary superfields.

On what can $f$ and $K$ depend? The function $f$ will be left-chiral if it
depends only on left-chiral elementary superfields $\Phi_n$, but not their
right-chiral complex conjugates. On the other hand, the superderivative of a
chiral superfield is not chiral, so we cannot freely include superderivatives
of the $\Phi_n$ in $f$. It is true that, acting on a superfield $S$ that is not
left-chiral (such as one involving complex conjugates of left-chiral
superfields), a pair of right superderivatives gives a left-chiral superfield,
because there are only two independent right superderivatives and they
anticommute:
\[
  \bar D_{\dot\alpha}
  \bigl(\bar D_{\dot\beta}\bar D_{\dot\gamma}S\bigr)=0\,.
\]
However, the $\mathcal{F}$-term of any function $f$ that is constructed in
this way makes a contribution to the action that is the same as that of a
$D$ term of some other composite superfield. Since the \(\bar D\)'s
anticommute, the most general left-chiral superfield formed by acting on a
general superfield \(S\) with two \(\bar D\)'s may be expressed in terms of
\(\bar D_{\dot\alpha}\bar D^{\dot\alpha}S\). If one of the left-chiral
superfields in the superpotential is of this form, then since each
\(\bar D_{\dot\alpha}\) annihilates
all the other superfields in the superpotential, we can write the whole
superpotential as
\(f=(\bar D_{\dot\alpha}\bar D^{\dot\alpha})h\) for some other
superfield $h$. Now,
\[
  \bar D_{\dot\alpha}\bar D^{\dot\alpha}
  (\bar\theta\bar\theta)=-4\,,
\]
With the integration order used below, the corresponding highest monomial is
\[
  \frac12\theta^2\bar\theta^2
  =-2\theta_1\theta_2
    \bar\theta_{\dot1}\bar\theta_{\dot2}.
\]
so, apart from spacetime derivatives that do not contribute to the action,
\((\bar D_{\dot\alpha}\bar D^{\dot\alpha})h\) is the coefficient of
\(-\bar\theta\bar\theta/4\) in \(h\). But, again apart from spacetime
derivatives, $[f]_{\mathcal{F}}$ is the coefficient of
\(\theta\theta\) in \(f\), so
  \([(\bar D_{\dot\alpha}\bar D^{\dot\alpha})h]_{\mathcal{F}}\) equals minus
  twice the \(D\)-component of \(h\), and therefore
\begin{equation}
  \int d^4x\,
  [(\bar D_{\dot\alpha}\bar D^{\dot\alpha})h]_{\mathcal{F}}
  =-2\int d^4x\,[h]_D\,.
  \tag{26.3.31}\label{eq:26.3.31}
\end{equation}
Thus we do not need to include terms in $f$ that depend on left-chiral
superfields of the form
\(\bar D_{\dot\beta}\bar D_{\dot\gamma}S\)---any such terms
will be included in the list of all possible $D$-terms. When $f$ is expressed
as a function only of elementary left-chiral superfields and not their
superderivatives or spacetime derivatives, it is known as the
\emph{superpotential}.

In contrast, the function $K$ is in general a real scalar function of both
left-chiral superfields $\Phi_n$ and their right-chiral complex conjugates
$\Phi_n^*$, as well as their superderivatives and spacetime derivatives,
known as the \emph{Kahler potential}. (Any right-chiral superfield is the
complex conjugate of a left-chiral superfield, so there is no loss of
generality in supposing $K$ to depend only on left-chiral superfields and
their complex conjugates.) However, not all $K$s obtained in this way will
yield distinct actions. For one thing, chiral superfields have no $D$-terms,
so two $K$s that differ by a chiral superfield make the same contribution to
the action.

It is also possible to change the form of $K$ without changing the action by a
partial integration in superspace. The $D$-term of \(D_\alpha S\), or its
conjugate, for an arbitrary superfield makes no contribution to the
action because
\begin{equation}
  \int d^4x\,[D_\alpha S]_D
  =\int d^4x\,[\bar D_{\dot\alpha}S]_D=0\,.
  \tag{26.3.32}\label{eq:26.3.32}
\end{equation}
Using the definition \eqref{eq:26.2.26}, the two covariant derivatives act
explicitly as
\[
  D_\alpha S
  =-\frac{\partial S}{\partial\theta^\alpha}
    +i\sigma^\mu_{\alpha\dot\alpha}\bar\theta^{\dot\alpha}
      \partial_\mu S,
  \qquad
  \bar D_{\dot\alpha}S
  =\frac{\partial S}{\partial\bar\theta^{\dot\alpha}}
    -i\theta^\alpha\sigma^\mu_{\alpha\dot\alpha}\partial_\mu S.
\]
Since \(S\) is a polynomial at most of second order in each of
\(\theta,\bar\theta\), the fermionic-coordinate derivative in
\(D_\alpha S\) is at most cubic in the combined coordinates and
therefore cannot have a non-zero $D$-term that is not a derivative, while the
remaining term in \(D_\alpha\) is a spacetime derivative, so its \(D\)-term is
also a
spacetime derivative, and therefore neither the first nor the second terms in
\(D_\alpha S\) can contribute to the integral in
Eq.~\eqref{eq:26.3.32}. Also, the superderivative acts distributively, so it
follows from Eq.~\eqref{eq:26.3.32} that we can integrate by parts in
superspace: for any two bosonic superfields $S_1$ and $S_2$,
\begin{equation}
  \int d^4x\,[S_1D_\alpha S_2]_D
  =-\int d^4x\,[S_2D_\alpha S_1]_D\,,
  \tag{26.3.33}\label{eq:26.3.33}
\end{equation}
In Sections 26.4 and 26.8 we will consider in detail the case where $f$ and
$K$ depend only on elementary superfields, but not their superderivatives or
ordinary derivatives.

We saw in the previous section that in theories in which parity is conserved,
the effect of the space inversion operator on a general scalar superfield is
to subject its arguments to the transformations
$x^\mu\rightarrow(\Lambda_P)^\mu{}_\nu x^\nu$ and
\((\theta,\bar\theta)\rightarrow(-i\bar\theta,-i\theta)\), with spinor
indices connected by \(\sigma^0\), and perhaps to multiply the superfield by
a phase \(\eta\). Under these transformations, the arguments \(x_\pm^\mu\) in
Eqs.~\eqref{eq:26.3.21} and \eqref{eq:26.3.22} are changed by
\begin{equation}
  x_\pm^\mu\longrightarrow
  (\Lambda_Px_\mp)^\mu\,,
  \tag{26.3.34}\label{eq:26.3.34}
\end{equation}
while \(\theta\rightarrow-i\sigma^0\bar\theta\) and
\(\bar\theta\rightarrow-i\bar\sigma^0\theta\). Thus space inversion takes left-chiral
superfields into right-chiral superfields, and vice-versa. The only
right-handed scalar superfield whose component fields involve creation and
annihilation operators for the same particles that are created and destroyed
by a left-handed scalar superfield $\Phi$ is
$\widetilde{\Phi}\mathrel{\propto}\Phi^*$, so $P^{-1}\Phi P$ must be
proportional to $\Phi^*$. By a suitable choice of phase of $\Phi$, we can
arrange that this transformation rule reads
\begin{equation}
  P^{-1}\Phi(x,\theta,\bar\theta)P
  =\Phi^\dagger(\Lambda_Px,-i\bar\theta,-i\theta)\,.
  \tag{26.3.35}\label{eq:26.3.35}
\end{equation}
In terms of component fields, this transformation is
\begin{equation}
  \begin{aligned}
  P^{-1}\phi(x)P&=\phi^*(\Lambda_Px),\\
  P^{-1}\psi_\alpha(x)P
    &=-i(\sigma^0)_{\alpha\dot\alpha}
      \bar\psi^{\dot\alpha}(\Lambda_Px),\\
  P^{-1}\mathcal{F}(x)P&=\mathcal{F}^*(\Lambda_Px).
  \end{aligned}
  \tag{26.3.36}\label{eq:26.3.36}
\end{equation}

There is another type of possible symmetry, known as $R$-symmetry, that is
important in some models of spontaneous supersymmetry breaking discussed in
Section 26.5, and that will also be used in proving no-renormalization
theorems in Section 27.6. As mentioned in Section 25.2, in theories of simple
$N=1$ supersymmetry an $R$-symmetry is invariance under a $U(1)$
transformation under which the left-handed components of the supersymmetry
generator \(Q_\alpha\) carry a non-vanishing quantum number, say \(-1\), in
which case their adjoints \(\bar Q_{\dot\alpha}\) carry \(+1\). Inspection of
Eq.~\eqref{eq:26.2.2} shows that the superspace coordinates have a
non-trivial transformation property under $R$-symmetry transformations:
\(\theta\) carries \(R\)-quantum number \(+1\), while \(\bar\theta\)
carries \(-1\). In addition,
the whole superfield may be given an $R$ quantum number. If we give a
left-chiral superfield $\Phi$ the $R$ quantum number $R_\Phi$, then its scalar
component $\phi$ has the same $R$ quantum number, while the left-handed spinor
component \(\psi_\alpha\) has \(R_\psi=R_\Phi-1\), and the auxiliary field
$\mathcal{F}$ has $R_{\mathcal{F}}=R_\Phi-2$. In particular, in order for the
superpotential term $\int d^4x\,[f]_{\mathcal{F}}$ to conserve $R$, the
superpotential itself must have $R_f=+2$, so if $f$ depends on a single
left-chiral superfield $\Phi$, then it must be proportional to
$\Phi^{2/R_\Phi}$. To put this another way, if $f(\Phi)$ is a pure mass term
proportional to $\Phi^2$ then we must choose $R_\Phi=+1$, while if $f(\Phi)$
is a pure interaction term proportional to $\Phi^3$, then we must choose
$R_\Phi=2/3$. On the other hand, inspection of Eq.~\eqref{eq:26.2.10} shows
that the $D$-term of a superfield has the same $R$ value as the superfield, so
in order for the term $\int d^4x\,[g]_D$ in the action to conserve $R$ it is
only necessary that $K$ have $R=0$, which will be the case if each term in
$K$ contains equal numbers of $\Phi$ and $\Phi^*$ factors, whatever $R$ value
we give to $\Phi$. Of course, there is no general reason why $R$-symmetry
should be respected by the action, or why it should not be spontaneously
broken.

\begin{center}
\(\ast\qquad\ast\qquad\ast\)
\end{center}

There are other ways of constraining superfields, to yield other types of
supermultiplets of fields. Among the more common are the \emph{linear}
superfields. To learn the conditions defining this sort of superfield, we note
that if $S$ is a general superfield, then we can form a chiral superfield
\begin{equation}
  S'\equiv\frac14
  \left(D^\alpha D_\alpha
    +\bar D_{\dot\alpha}\bar D^{\dot\alpha}\right)S\,.
  \tag{26.3.37}\label{eq:26.3.37}
\end{equation}
This is a chiral superfield because it is the sum of
\(\frac14D^\alpha D_\alpha S\), which is right-chiral, and
\(\frac14\bar D_{\dot\alpha}\bar D^{\dot\alpha}S\), which is left-chiral.
Its components are
given in terms of the components of $S$ by
\begin{equation}
  C'=N\,,
  \tag{26.3.38}\label{eq:26.3.38}
\end{equation}
\begin{equation}
  \omega'_\alpha
  =\lambda_\alpha-i\sigma^\mu_{\alpha\dot\alpha}
    \partial_\mu\bar\omega^{\dot\alpha}\,,
  \tag{26.3.39}\label{eq:26.3.39}
\end{equation}
\begin{equation}
  M'=-\partial_\mu V^\mu\,,
  \tag{26.3.40}\label{eq:26.3.40}
\end{equation}
\begin{equation}
  N'=D+\Box C\,,
  \tag{26.3.41}\label{eq:26.3.41}
\end{equation}
\begin{equation}
  V_\mu'=-\partial_\mu M\,,
  \tag{26.3.42}\label{eq:26.3.42}
\end{equation}
\begin{equation}
  \lambda'=D'=0\,.
  \tag{26.3.43}\label{eq:26.3.43}
\end{equation}
A multiplet $S$ is said to be linear if the superfield $S'$ defined in this
way vanishes
\begin{equation}
  \left(D^\alpha D_\alpha
    +\bar D_{\dot\alpha}\bar D^{\dot\alpha}\right)S=0\,,
  \tag{26.3.44}\label{eq:26.3.44}
\end{equation}
or in terms of its components,
\begin{equation}
  N=M=\partial_\mu V^\mu=0,\qquad
  \lambda_\alpha
    =i\sigma^\mu_{\alpha\dot\alpha}\partial_\mu
      \bar\omega^{\dot\alpha},\qquad
  D=-\Box C\,.
  \tag{26.3.45}\label{eq:26.3.45}
\end{equation}
This leaves four independent bosonic fields, $C$ and the three components of
$V_\mu$ subject to the condition $\partial_\mu V^\mu=0$, and four independent
fermionic fields, the two complex components of \(\omega_\alpha\) and their
conjugates. We
will see in Section 26.6 that the current superfields, whose $V_\mu$-terms are
the conserved currents associated with symmetry transformations, are linear
superfields.
